\documentclass[../main]{subfiles}
\begin{document}

\chapter{THREAT MODEL AND PROBLEM FORMULATION}
\label{ch:threat_model}

\section{System Model: The Federated Learning System}
\label{sec:sys_model}

% TODO: N clients, central server as aggregator,
% communication graph, partial participation model.

\section{Adversary Model}
\label{sec:adv_model}

% TODO: Adversary capabilities (knowledge, persistence, coordination bound f < n/2).

\subsection{Inner Product Manipulation (IPM) Attacks}
\label{sec:ipm}

% TODO: Formal definition of IPM. Adversary constructs g_byz such that
% <g_byz, g_honest> < 0. Why cosine screening alone is insufficient.

\subsection{Sybil Attacks}
\label{sec:sybil}

% TODO: Single physical adversary controls multiple virtual identities.
% Effective Byzantine fraction inflated beyond f.

\subsection{Sign-Flip Attacks}
\label{sec:sign_flip}

% TODO: g_byz = -alpha * g_honest. Analysis of damage under FedAvg.

\subsection{Label-Flipping and Backdoor Attacks}
\label{sec:label_flip}

% TODO: Data-poisoning variant; how Aegis provides incidental robustness.

\section{Problem Statement and Design Goals}
\label{sec:problem_stmt}

% TODO: Formal objectives: (i) convergence guarantee for honest clients,
% (ii) O(kd) aggregation complexity, (iii) no trusted server gradient.

\begin{problem}[Byzantine-Resilient Aggregation]
\label{prob:main}
% TODO: Formal problem statement.
\end{problem}

\section{Formal Threat Bounds}
\label{sec:threat_bounds}

% TODO: Theorem stating max tolerable Byzantine fraction and conditions.

\begin{assumption}[Gradient Boundedness]
\label{asm:grad_bound}
% TODO
\end{assumption}

\begin{assumption}[Honest Majority]
\label{asm:honest_maj}
% TODO
\end{assumption}

\end{document}
